\section{Motivation and Challenges}\label{sec:motivation_challenges}

\begin{figure}[tp!]
\centering
  \includegraphics[width=\linewidth]{figures/motivation/traces.png}
\vspace{-15pt}
\caption{Availability of A100 GPUs in two zones in Google Cloud over 8-hr period. We request 8 GPUs in each zone. The trace was collected in April 2024.} 
\label{fig:trace}
\vspace{-0.5cm}
\end{figure}

\subsection{Why use heterogeneous, geo-distributed GPUs?}\label{motivation}

\textbf{Homogeneous high-end GPU clusters are scarce.} The widespread adoption of ML, and hardware vendors' inability to keep up with this pace, has significantly increased GPU demand. Several studies report limited GPU availability across public cloud providers~\cite{Zongheng23Skypilot, Cheng23availability, chugh23availability, strati24crossregion, Thorpe23Bamboo}. 
For example, Figure \ref{fig:trace} plots the availability of on-demand A100 GPUs in two zones in GCP. In one zone, it took 7 hours to allocate 8 A100 GPUs, while in the other zone, the requested number of GPUs was not attained within the 8-hour window.
Our findings align with the AWS GPU availability plot from Guo et al.~\cite{guo2024cephalo}, which shows that high-end GPUs such as A100 and H100 are difficult to aquire, while mid-tier GPUs (A10G, V100, T4) have higher but still limited and dynamic availability. %\qh{Better have a trace plot here for clearer demonstration. I remember we have such data}
Using GPUs across zones with heterogeneous types gives ML jobs the opportunity to use more GPUs to further increase throughput, as shown in Figure~\ref{fig:motivation_het}. 


\textbf{Power and cooling limits \#GPUs per datacenter.}  From 2020 to 2025, the size of state-of-the-art ML models has increased by roughly 1200$\times$~\cite{gpt-2, gpt-4, llama-4}, while per-GPU memory capacity has only increased by 8$\times$~\cite{nvidia-list}. This requires hyperscalers to deploy more and more GPUs~\cite{ms-newAIdatacenter, meta-newAIdatacenter}. 
However, the high power and cooling requirements of high-end GPUs limits the amount that can be deployed per datacenter~\cite{tapas}. The next generation of ML models may need to train on GPUs across multiple availability zones~\cite{geodistributedblog1, geodistributedblog2, heatchallengeblog, palak2024improvingtrainingtimegpu}.

\textbf{Using old GPUs can reduce embodied carbon.} Embodied carbon (greenhouse gas emissions associated with manufacturing to disposal) is a major source of datacenter  emissions~\cite{greenSKU, carbonExplorer, metaEmbodiedCarbon}.
Although users prefer to allocate the latest GPUs for their ML jobs, older GPUs are abundant as the typical lifetime for ML servers in hyperscaler datacenters is $\sim$6 years~\cite{schneider2025lifecycleGoogleAI, greenSKU}. Finding optimal ways to spread jobs across heterogeneous GPUs will enable leveraging older GPUs for longer to better amortize embodied carbon. 


\subsection{Challenges with Heterogeneous ML Clusters}\label{challenges}

\textbf{C1: Quickly searching a vast configuration space.} 
Considering heterogeneous and geo-distributed GPUs creates a vast and complex search space.
The ML developer needs to decide how many GPUs to use and how to group them across VMs.
This complexity increases further, when accounting for job parallelization plans within each allocation. Furthermore, the optimal allocation and partitioning strategies depend on user objectives and constraints.
A planner needs to \textit{quickly} navigate the large configuration space to adapt to dynamic resource availability in cloud and on-premise environments~\cite{Thorpe23Bamboo, guo2024cephalo, weng22mlaas, Gu19Tiresias}, since maximizing throughput requires adapting parallelization strategies with changes to cluster topology~\cite{Athlur22Varuna, wagenlander24tenplex}. 


Table~\ref{table:planners} shows that Metis~\cite{Um24Metis}, FlashFlex~\cite{yan2024flashflex}, Cephalo~\cite{guo2024cephalo}, Atlas~\cite{palak2024atlas}, and DTFM~\cite{yuan2023decentralized}  explore parts of this large search space. Atlas~\cite{palak2024atlas} and DTFM~\cite{yuan2023decentralized}  target geo-distributed training, but do not consider heterogeneous GPU types, and do not decide the various parallelism degrees: instead, they take as input the parallelism degrees, and assign these degrees in the available zones.
On the other hand, Metis~\cite{Um24Metis}, FlashFlex~\cite{yan2024flashflex} and Cephalo~\cite{guo2024cephalo},
consider heterogeneous GPU types, but overlook geo-distributed training, and are quite inefficient for dynamic environments. Metis needs a few hours to devise a plan for a 16-GPU cluster (A100 and V100)\footnote{with max\_permutation\_length and device group variance set to 10 and 0.5 respectively, according to the paper~\cite{Um24Metis}}, making frequent reevaluation infeasible as GPU availability changes. Cephalo~\cite{guo2024cephalo} has a search time of ~300 sec on a cluster of 64 GPUs\footnote{reported on the paper~\cite{guo2024cephalo}}, but it is limited only to Fully Sharded Data Parallelism. FlashFlex~\cite{yan2024flashflex} has a short runtime, but provides inaccurate runtime estimations, leading to suboptimal plans. 
Furthermore, these planners only optimize for \textit{throughput}, ignoring budget constraints, and cost (dollars per iteration), which affect the optimal configuration (\S\ref{sec:eval_constraints}). 


\textbf{C2: Accurately simulating memory footprint and iteration time.} Most planners use analytical models or simulations to evaluate a configuration (parallelism degrees, microbatch sizes, etc) on a given cluster setup, since it is impractical and very expensive to deploy and profile every configuration. This evaluation usually includes two stages: 1) memory footprint estimation, to identify whether a configuration is valid (i.e., it does not lead to OOM errors) and 2) iteration time estimation, to determine performance. 

{\setlength{\textfloatsep}{-0cm}
\begin{figure}[tp!]
\centering
\vspace{5pt}
  \includegraphics[width=\columnwidth]{figures/motivation/opt350_mem.pdf}%
\caption{Peak Memory estimations of various baselines compared to the actual peak memory, for the OPT-350M model on a homogeneous cluster of 4 Grace-Hopper per node.  $N$ stands for the number of nodes,  $gbs$ is the global batch size, $pp$ pipeline parallelism, $tp$ tensor parallelism, $dp$ data parallelism, and $mbs$ the microbatch size. }
\label{fig:motivation_sim_mem}
%\vspace{-0.75cm}
\end{figure}}


Unfortunately, these estimations are often inaccurate, resulting in suboptimal or invalid plans. First,  planners either completely ignore memory footprint~\cite{li2022amp}, or underestimate the amount of memory requirements during training~\cite{Athlur22Varuna, Jang2023Oobleck}, omitting activations, optimizer states, and memory fragmentation, or assume the training memory footprint is uniform across all devices and pipeline stages~\cite{yan2024flashflex, Lin24nnScaler}. As a result, these systems may find plans that cause OOM errors, when deployed. Figure \ref{fig:motivation_sim_mem} compares memory footprint estimations with the real footprint on a homogeneous cluster of up to 16 Grace-Hopper nodes for the OPT-350M model, showing that planners can be 25-95\% off when estimating memory footprint.  Second, planners often poorly model training time, due to incorrect assumptions about network bandwidth and ignoring communication-computation overlap~\cite{Zhang17Poseidon, megatron-deepspeed}.

Modeling memory footprint and iteration time becomes even more complex with heterogeneity~\cite{Mo24Heet}. Different GPU generations vary in compute performance (e.g., TFLOPs), and memory capacity~\cite{zhu2024nanoflow}. Thus, a configuration that fits in one GPU, might cause OOM errors in another GPU. Additionally, stragglers and network bandwidth differences (especially in geo-distributed setups~\cite{strati24crossregion}) must be considered for accurate timing estimations. Planners must accurately model compute, memory, and network bandwidth heterogeneity - yet, as shown in Table~\ref{table:planners}, most systems overlook this.


\textbf{C3: Supporting both heterogeneous plans and seamless elasticity in a real distributed training framework.}   
Most training frameworks, i.e., systems that train a model on a set of devices, assume homogeneous clusters and job configuration plans. For example, widely used and high-performing frameworks such as Megatron~\cite{shoeybi2020megatronlm} and DeepSpeed~\cite{megatron-deepspeed} assume \textit{uniform} parallelism degrees for the entire training job (e.g. DP=2, PP=6, TP=1). This limits efficiency in heterogeneous configurations, where various parallelism degrees per training subproblem (e.g., using PP=6, with first 3 stages having TP=4, and the next 3 stages having TP=2) help load-balance compute and memory on GPU nodes with different resources. Thus, a framework should accommodate heterogeneous degrees of parallelism. Furthermore, the framework should seamlessly reconfigure and adapt to dynamic GPU availability (Figure \ref{fig:trace}), as long job reconfiguration times in response to resource changes are wasteful~\cite{Thorpe23Bamboo, wagenlander24tenplex}. Although related works propose elastic systems~\cite{Athlur22Varuna, Thorpe23Bamboo, Duan24Parcae} or systems that support heterogeneity~\cite{Miao23SDPipe, yan2024flashflex}, there is no open-source system that supports both.

